%% kapa: exam=1 round=36 year=2025 period=1 form=A subject=real_estate subject_ko=부동산학원론 seq=26 number=106
%% source: Q-Net 감정평가사 기출문제 (cst00309) artlSeq=5231525 file=raw/1008/5231525/2025년 제36회 감정평가사 1차시험 1교시 문제지 원본.pdf
%% conversion: kordoc 4.24.51 -> tools/gen1.py (auto); review=proofread-sonnet; flags=none
\begin{question}{106}
국토의 계획 및 이용에 관한 법률상 다음에 해당하는 계획은?
\begin{framedtext}
토지의 이용 및 건축물이나 그 밖의 시설의 용도ㆍ건폐율ㆍ용적률ㆍ높이 등을 완화하는 용도구역의 효율적이고 계획적인 관리를 위하여 수립하는 계획을 말한다.
\end{framedtext}
\begin{choices}
  \item 성장관리계획
  \item 도시혁신계획
  \item 복합용도계획
  \item 지구단위계획
  \item 공간재구조화계획
\end{choices}
\end{question}
\answer{5}
